% =========================================================
\chapter{Repositorio de código: guía de navegación}
\label{chap:anexo_codigo}
% =========================================================

El código fuente completo de este Trabajo Fin de Máster se encuentra
disponible de forma pública en el siguiente repositorio de GitHub:

\begin{center}
    \url{https://github.com/juanhdezz/TFM-Glucose-Prediction}
\end{center}

Una wiki interactiva generada automáticamente a partir del código, que
documenta en profundidad cada módulo y subsistema, está disponible en:

\begin{center}
    \url{https://deepwiki.com/juanhdezz/TFM-Glucose-Prediction}
\end{center}

Este apéndice actúa como guía de navegación del repositorio: describe la
organización de directorios, el propósito de cada subsistema y el orden
lógico en que deben consultarse los scripts para reproducir los
experimentos o extender el trabajo.

% ---------------------------------------------------------
\section{Estructura general del repositorio}
\label{sec:anexo_estructura}
% ---------------------------------------------------------

El repositorio se organiza en seis subsistemas principales, cada uno con
responsabilidades bien delimitadas:

\begin{description}
    \item[\texttt{src/}] Pipeline experimental principal. Contiene los
    scripts numerados que ejecutan, en orden, la generación de ventanas,
    la partición por paciente, el balanceo demográfico, el entrenamiento
    de los modelos LSTM y la evaluación de predicciones.

    \item[\texttt{EDA/}] Análisis Exploratorio de Datos. Agrupa los
    scripts Python y \textit{notebooks} Jupyter que generan las figuras
    descriptivas de los tres \textit{datasets}, los \textit{heatmaps}
    cruzados y los informes de desequilibrio por \textit{fold}.

    \item[\texttt{analysis\_relief/}] Marco de análisis de resultados
    \textit{post-hoc}. Implementa la carga de resultados, las pruebas
    estadísticas no paramétricas (Friedman, Nemenyi), las
    visualizaciones de publicación y la generación automática de tablas
    \LaTeX.

    \item[\texttt{folds\_analysis/}] Informes de desequilibrio por
    \textit{fold} y \textit{dataset}, en formato texto, generados como
    \textit{proxy} del ratio de desbalanceo dinámico inter-\textit{fold}.

    \item[\texttt{batch\_scripts/}] Ficheros de cola SLURM para la
    ejecución en el clúster HPC. Automatizan el envío de trabajos de
    balanceo y entrenamiento a gran escala.

    \item[\texttt{memoria/}] Documento \LaTeX\ del TFM. Contiene los
    capítulos fuente (\texttt{.tex}), las figuras referenciadas y el
    PDF compilado.
\end{description}

Los datos clínicos en formato Parquet (\texttt{data/}) y los artefactos
de resultados (\texttt{results/}, \texttt{cache/}) están excluidos del
control de versiones mediante \texttt{.gitignore}, dado su tamaño y
carácter sensible.

% ---------------------------------------------------------
\section{Configuración del entorno de desarrollo}
\label{sec:anexo_entorno}
% ---------------------------------------------------------

El proyecto requiere Python~3.10 o superior. Todas las dependencias
están fijadas en \texttt{requirements.txt} para garantizar la
reproducibilidad exacta de los experimentos. Los grupos de librerías
principales son:

\begin{itemize}
    \item \textbf{Manipulación y almacenamiento de datos}:
    \texttt{pandas==2.2.3}, \texttt{numpy==2.2.3}, \texttt{pyarrow}
    (lectura/escritura de ficheros Parquet).

    \item \textbf{Aprendizaje automático y corrección de desequilibrio}:
    \texttt{scikit-learn==1.8.0}, \texttt{imbalanced-learn==0.14.1},
    \texttt{scipy==1.17.1}.

    \item \textbf{Visualización y \textit{notebooks}}:
    \texttt{matplotlib==3.10.8}, \texttt{seaborn==0.13.2},
    \texttt{ipykernel}.
\end{itemize}

Para inicializar el entorno local se recomienda crear un entorno virtual
e instalar las dependencias:

\begin{verbatim}
python -m venv .venv
source .venv/bin/activate          # Linux/macOS
pip install -r requirements.txt
\end{verbatim}

\subsection{Organización del directorio de datos}
\label{subsec:anexo_datos}

Los tres \textit{datasets} clínicos deben ubicarse bajo \texttt{data/}
siguiendo la estructura esperada por los scripts de carga:

\begin{verbatim}
data/
+-- T1DiabetesGranada/
|   +-- Patient_info.csv
|   +-- Glucose_measurements.parquet
+-- REPLACE-BG/
|   +-- Glucose_measurements_REPLACE-BG_FILTERED.parquet
+-- DIATREND/
    +-- Glucose_measurements_DiaTrend_FILTERED.parquet
\end{verbatim}

La resolución de rutas se gestiona programáticamente en cada script
mediante utilidades de búsqueda del directorio raíz del workspace, lo
que hace el código portable independientemente de la ruta absoluta de
instalación.

\subsection{Ejecución local frente a clúster HPC}
\label{subsec:anexo_hpc}

El pipeline soporta dos modalidades de ejecución:

\begin{enumerate}
    \item \textbf{Ejecución local}: adecuada para el análisis exploratorio
    (\texttt{EDA/}), el \textit{testing} ligero de scripts y la
    compilación del documento \LaTeX. Se ejecutan las sesiones Jupyter
    interactivas o los scripts directamente desde la línea de comandos.

    \item \textbf{Ejecución en clúster HPC}: obligatoria para las tareas
    computacionalmente intensivas: generación de ventanas deslizantes
    (\texttt{src/1-generate\_windows.py}), balanceo demográfico con las
    diez técnicas (\texttt{src/2-server\_balancing.py}) y entrenamiento
    de los modelos LSTM (\texttt{src/3-train.py}). Los trabajos se
    gestionan mediante ficheros SLURM ubicados en \texttt{batch\_scripts/},
    que dirigen los artefactos de salida hacia \texttt{data/output/}.
\end{enumerate}

% ---------------------------------------------------------
\section{Pipeline experimental (\texttt{src/})}
\label{sec:anexo_pipeline}
% ---------------------------------------------------------

El directorio \texttt{src/} implementa el flujo experimental
\textit{end-to-end}: desde los datos clínicos en bruto hasta los
modelos evaluados. Los scripts están numerados para reflejar el orden de
ejecución obligatorio.

\subsection{Paso 1: Generación de ventanas deslizantes
            (\texttt{1-generate\_windows.py})}
\label{subsec:anexo_ventanas}

Transforma las series temporales CGM en instancias de aprendizaje
supervisado mediante ventanas deslizantes. El parámetro de longitud de
historia se fija en $H=8$ muestras (equivalente a 40~minutos con
frecuencia de muestreo de 5~min) y el horizonte de predicción en
$PH=4$ (20~minutos), bajo restricciones estrictas de continuidad
temporal: ninguna ventana puede cruzar lagunas de datos superiores al
umbral definido. Cada ventana hereda los metadatos demográficos del
paciente al que pertenece (\texttt{patient\_id}, \texttt{age\_group},
\texttt{sex}).

\subsection{Paso 2a: Partición Group K-Fold
            (\texttt{2-folds\_creation.py})}
\label{subsec:anexo_folds}

Divide el conjunto de ventanas en 5~\textit{folds} usando
\texttt{GroupKFold} de \textit{scikit-learn}, con el paciente como
unidad indivisible de agrupación. El script añade columnas
\texttt{fold\_*} al Parquet de ventanas para que cada fila quede
asignada a exactamente un \textit{fold} de test. Esta estrategia
garantiza la separación completa entre pacientes en entrenamiento y
evaluación, eliminando la fuga de información inherente a los esquemas
de validación cruzada aleatorios sobre señales fisiológicas continuas.

\subsection{Paso 2b: Balanceo demográfico
            (\texttt{2-server\_balancing.py})}
\label{subsec:anexo_balanceo}

Aplica las diez técnicas de corrección de desequilibrio demográfico
descritas en la Sección~\ref{sec:pipeline_entrenamiento} del cuerpo
principal de la memoria. El balanceo se ejecuta \textbf{exclusivamente}
sobre la partición de entrenamiento de cada \textit{fold}, dejando
intactas las particiones de validación y test para reflejar la
distribución demográfica real. Los \textit{datasets} balanceados se
serializan en Parquet con una nomenclatura que codifica el
\textit{dataset}, la dimensión demográfica (\texttt{age} o \texttt{sex})
y la técnica aplicada, por ejemplo:

\begin{verbatim}
T1DiabetesGranada_age_smote_fold0.parquet
REPLACE-BG_sex_undersampling_fold3.parquet
\end{verbatim}

El módulo auxiliar \texttt{src/balancing\_viz.py} genera figuras de
diagnóstico que muestran la distribución demográfica antes y después de
cada técnica.

\subsection{Paso 3: Entrenamiento LSTM (\texttt{3-train.py})}
\label{subsec:anexo_entrenamiento}

Descubre automáticamente los ficheros Parquet balanceados generados en
el paso anterior y entrena una red neuronal recurrente LSTM para cada
combinación de (\textit{dataset}, dimensión, técnica, \textit{fold}).
Las funciones de pérdida están definidas en \texttt{0-loss\_functions.py}
e incluyen RMSE estándar y pérdidas clínicas penalizadas que castigan de
forma asimétrica los errores en zonas de hipoglucemia severa (TBR-2) e
hiperglucemia severa (TAR-2). La reproducibilidad se garantiza mediante
semillas aleatorias fijas en todos los experimentos.

\subsection{Paso 4: Evaluación e inferencia
            (\texttt{4-test\_predictions.py})}
\label{subsec:anexo_evaluacion}

Ejecuta la inferencia sobre las particiones de test y calcula métricas
desagregadas por rango glucémico (TBR-2, TBR-1, TIR, TAR-1, TAR-2):
RMSE, MAE, MAPE y distribución de zonas en la Clarke Error Grid. Los
resultados se compilan en ficheros de informe plano (formato MTA) y se
agregan entre \textit{folds} para su análisis estadístico posterior.

% ---------------------------------------------------------
\section{Análisis exploratorio de datos (\texttt{EDA/})}
\label{sec:anexo_eda}
% ---------------------------------------------------------

El directorio \texttt{EDA/} contiene \textit{notebooks} Jupyter y
scripts Python organizados en dos niveles:

\subsection{Inspección por \textit{dataset}
            (\texttt{notebooks/01\_EDA/})}
\label{subsec:anexo_eda_individual}

Cada \textit{dataset} dispone de su propio \textit{notebook} de
inspección inicial:

\begin{itemize}
    \item \texttt{01\_Inspection\_T1DiabetesGranada.ipynb} /
          \texttt{\_v2.ipynb}
    \item \texttt{01\_Inspection\_ReplaceBG.ipynb}
    \item \texttt{01\_Inspection\_DIATREND.ipynb}
\end{itemize}

Estos \textit{notebooks} verifican la cadencia de muestreo, detectan
patrones MNAR (\textit{Missing Not At Random}), caracterizan la
distribución demográfica y calculan métricas clínicas de referencia
(TIR/TBR/TAR) sin balanceo.

\subsection{Análisis cruzado entre \textit{datasets}
            (\texttt{EDA/cross/})}
\label{subsec:anexo_eda_cross}

Los scripts de análisis cruzado generan las figuras comparativas que
aparecen en el cuerpo del TFM: distribuciones de sexo y edad por
\textit{dataset}, \textit{heatmaps} edad$\times$sexo, y diagramas de
desequilibrio por \textit{fold}. Los artefactos se nombran con el
prefijo \texttt{B1*} y se referencian directamente desde el documento
\LaTeX.

El subdirectorio \texttt{folds\_analysis/} contiene informes en formato
texto con los ratios de desequilibrio calculados para cada técnica de
balanceo, \textit{dataset} y \textit{fold}, organizados por
\textit{dataset}:

\begin{verbatim}
folds_analysis/
+-- DIATREND/
|   +-- diatrend_age_jittering.txt
|   +-- diatrend_age_smote.txt
|   +-- ...
+-- REPLACE-BG/
+-- T1DiabetesGranada/
\end{verbatim}

% ---------------------------------------------------------
\section{Marco de análisis de resultados
         (\texttt{analysis\_relief/})}
\label{sec:anexo_analisis}
% ---------------------------------------------------------

El paquete \texttt{analysis\_relief/} es el subsistema de análisis
\textit{post-hoc}. Su función es agregar los resultados experimentales
generados en el clúster, ejecutar las pruebas estadísticas y producir
las figuras y tablas que se incorporan al capítulo de resultados de la
memoria.

\subsection{Configuración y carga de datos
            (\texttt{analysis/config.py}, \texttt{loader.py})}
\label{subsec:anexo_config}

El módulo \texttt{config.py} define los nombres canónicos de
\textit{datasets}, dimensiones demográficas, técnicas de balanceo y
sus ordenaciones estándar. \texttt{loader.py} ingesta los Parquet de
resultados distribuidos generados por el HPC y los unifica en
estructuras \textit{pandas} limpias y alineadas, listas para el
análisis.

\subsection{Validación estadística (\texttt{stats.py})}
\label{subsec:anexo_stats}

Implementa la secuencia completa de validación estadística no
paramétrica:

\begin{enumerate}
    \item \textbf{Test de Friedman}: contraste global de que al menos
    una técnica difiere significativamente del resto en cada combinación
    (\textit{dataset}, dimensión, métrica).
    \item \textbf{Test \textit{post-hoc} de Nemenyi}: identifica qué
    pares de técnicas difieren significativamente, con corrección por
    comparaciones múltiples. Los resultados se visualizan como diagramas
    de diferencia crítica (CD).
    \item \textbf{Recuento de Borda}: agrega los rankings de las
    técnicas a través de \textit{folds} y métricas para obtener un
    ranking global robusto.
    \item \textbf{Consistencia de Spearman}: mide la estabilidad del
    ranking entre \textit{folds} de validación cruzada.
\end{enumerate}

\subsection{Módulos de visualización (\texttt{viz/})}
\label{subsec:anexo_viz}

El directorio \texttt{viz/} contiene módulos de \textit{plotting}
especializados que generan el catálogo completo de figuras de la tesis:
distribuciones de RMSE por técnica, diagramas CD de Nemenyi,
\textit{bump charts} de estabilidad de ranking entre \textit{folds},
\textit{lollipop charts} de Borda, y gráficos de tamaño de conjunto de
entrenamiento por técnica.

\subsection{Salidas clínicas (\texttt{outputs\_clinical/})}
\label{subsec:anexo_clinical}

Traduce las métricas predictivas en interpretaciones clínicamente
relevantes: gráficos de error glucémico diferencial ($\Delta$RMSE
respecto al baseline sin balanceo), \textit{heatmaps} de balanceo
segmentados por edad y sexo, gráficos de compromiso TIR/TBR y análisis
de los cambios en el tamaño del conjunto de entrenamiento inducidos por
cada técnica de remuestreo.

% ---------------------------------------------------------
\section{Documento de la memoria (\texttt{memoria/})}
\label{sec:anexo_memoria}
% ---------------------------------------------------------

El directorio \texttt{memoria/} contiene el documento \LaTeX\ completo
del TFM. El fichero raíz es \texttt{tfm.tex}, que importa los capítulos
desde \texttt{capitulos/}:

\begin{verbatim}
memoria/
+-- tfm.tex
+-- capitulos/
|   +-- 01_Introduccion.tex
|   +-- 02_Estado_del_Arte.tex
|   +-- 03_Desarrollo_del_TFM.tex
|   +-- 04_Resultados_y_Discusion.tex
|   +-- 05_Conclusiones.tex
+-- EDA/
+-- tfm.pdf
\end{verbatim}

Los ficheros auxiliares de compilación (\texttt{*.aux}, \texttt{*.toc},
\texttt{*.log}, \texttt{*.bbl}, \texttt{*.synctex.gz}) están excluidos
del repositorio. Para compilar el documento se recomienda usar
\texttt{latexmk}:

\begin{verbatim}
cd memoria/
latexmk -pdf -interaction=nonstopmode tfm.tex
\end{verbatim}